\documentclass{scrartcl}

% encoding:
\usepackage[utf8]{inputenc}
\usepackage[T1]{fontenc}

% fonts and math packages:
\usepackage{amsthm}
\usepackage[fleqn, leqno]{amsmath}
\setlength\mathindent{1.6cm}
\everydisplay{\displayindent=1cm}
\usepackage{amsfonts}
\usepackage{amssymb}
% main font:
\usepackage{newtxtext}
\usepackage{newtxmath}
\usepackage{textcomp}

% links:
\usepackage{xcolor,colortbl,graphicx}
\definecolor{linkcol}{rgb}{0,0.1,0.4}
\definecolor{citecol}{rgb}{0,0.2,0.2}
\usepackage[colorlinks=true,
            linkcolor=linkcol,
            urlcolor=linkcol,
            citecolor=citecol]{hyperref}

% for enumerate with roman and alphabetical labels:
\usepackage{enumerate}

% additional packages:
\usepackage{fancyvrb,paralist}
\usepackage{multirow,rotating}
\usepackage{eqlist}
\usepackage{mdwlist}

\usepackage[english]{babel} % English language/hyphenation
\usepackage{graphicx} % for including pictures
\usepackage{tikz}

% nicer tables:
\usepackage{booktabs}
\renewcommand{\arraystretch}{1.1} % more space between table rows

% general layout:
\usepackage{url,ellipsis}
\usepackage[final=true,step=1]{microtype}
\usepackage{ragged2e}

% section heading layout:
\setkomafont{section}{%
\renewcommand{\bfdefault}{sb}%
%\usefont{T1}{qhv}{b}{n}\selectfont
%\bfseries\Large
\fontfamily{put}\bfseries\Large%
\renewcommand{\bfdefault}{bx}%
}
\setkomafont{subsection}{%
\renewcommand{\bfdefault}{sb}%
\fontfamily{put}\bfseries\large%
\renewcommand{\bfdefault}{bx}%
}

% make box and diamond larger, and vdash (too small in newtxmath):
\usepackage{scalerel}
\let\oldbox\Box
\renewcommand\Box{\scaleobj{1.2}{\oldbox}}
\let\olddiamond\Diamond
\renewcommand\Diamond{\scaleobj{1.2}{\olddiamond}}
\let\oldvdash\vdash
\renewcommand\vdash{\scaleobj{1.3}{\oldvdash}}

% make \to smaller and with less spacing to the sides:
\let\oldto\to
\renewcommand\to{\scaleobj{0.95}{\,\oldto\,}}
\let\oldleftrightarrow\leftrightarrow
\renewcommand\leftrightarrow{\scaleobj{0.95}{\,\oldleftrightarrow\,}}

% define strictif and boxright:
\DeclareSymbolFont{symbolsC}{U}{txsyc}{m}{n}
\DeclareMathSymbol{\strictif}{\mathrel}{symbolsC}{74}
\DeclareMathSymbol{\boxright}{\mathrel}{symbolsC}{128}

% define boxes and diamonds for specialist logics:
\DeclareMathOperator{\Kn}{\mathsf{K}{}}
\DeclareMathOperator{\Ind}{\mathsf{Ind}{}}
\DeclareMathOperator{\Mi}{\mathsf{M}}
\DeclareMathOperator{\Po}{\mathsf{P}}
\DeclareMathOperator{\Bel}{\mathsf{B}}
\DeclareMathOperator{\EKn}{\mathsf{E}}
\DeclareMathOperator{\CKn}{\mathsf{C}}
\DeclareMathOperator{\tP}{\mathsf{P}}
\DeclareMathOperator{\tF}{\mathsf{F}}
\DeclareMathOperator{\tG}{\mathsf{G}}
\DeclareMathOperator{\tH}{\mathsf{H}}
\DeclareMathOperator{\tX}{\mathsf{X}}
\DeclareMathOperator{\tY}{\mathsf{Y}}
\DeclareMathOperator{\tS}{\mathsf{S}}
\DeclareMathOperator{\tU}{\mathsf{U}}
\DeclareMathOperator{\tN}{\mathsf{N}}
\DeclareMathOperator{\Ob}{\mathsf{O}}
\DeclareMathOperator{\Pe}{\mathsf{P}}
\DeclareMathOperator{\tA}{\mathsf{A}}
\DeclareMathOperator{\Always}{\mathsf{A}}
\DeclareMathOperator{\Sometimes}{\mathsf{S}}
\DeclareMathOperator{\Mostly}{\mathsf{M}}

% better smallcaps:
\newcommand{\smallcaps}[1]{{\text{\footnotesize \MakeUppercase{#1}}}}
\newcommand{\stit}{\mathop{\textsf{\smallcaps{stit}}}}
\newcommand{\OK}{\mathop{\textsf{\smallcaps{OK}}}}

% Trees:
\usepackage{xyling}
\usepackage{xfp}
\usepackage{ifthen}
\definecolor{annotcol}{rgb}{0.3, 0.4, 0.5}
\definecolor{red}{rgb}{1,0,0}
\def\checkmark{\color{annotcol}\tikz\fill[scale=0.4](0,.35) -- (.25,0) -- (1,.7) -- (.25,.15) -- cycle;}
\newcommand{\xmark}{\ding{55}}%
\newcommand{\tree}[2][0]{\Treek[#1]{0.7}{#2}}
%\newcommand{\dotline}{\GB{d}{.}}
\newcommand{\wlabel}[2]{\Kk[#1]{0}{\ifthenelse { \equal {#2} {} } {} {\ensuremath{(#2)}}}}
\newcommand{\ticklabel}[1]{\Kk[#1]{0}{\checkmark}}
\newcommand{\annotlabel}[2]{\Kk[#1]{0}{{\color{annotcol}#2}}}
\newcommand{\llabel}[2]{\Kk[-#1]{0}{{\color{annotcol}#2}}}
\newcommand{\branchbelow}{\B{ddl}\B{ddr}}
\newcommand{\closed}{\Below{\sffamily x}}
\newcommand{\treenode}[5]{\llabel{\fpeval{(#1)-1}}{#2} #3 \wlabel{#1}{#4} \annotlabel{\fpeval{(#1)+10}}{#5}}
\newcommand{\treenodeticked}[5]{\llabel{\fpeval{(#1)-1}}{#2} #3 \wlabel{#1}{#4} \annotlabel{\fpeval{(#1)+10}}{#5}  \ticklabel{\fpeval{(#1)+18}}  }
\newcommand{\nnode}[5]{\treenode{#1}{#2}{\K{#3}}{#4}{#5}}
\newcommand{\nnodeticked}[5]{\treenodeticked{#1}{#2}{\K{#3}}{#4}{#5}}
\newcommand{\nnodeclosed}[5]{\treenode{#1}{#2}{\K{#3}\closed}{#4}{#5}}
\newcommand{\dotbelowbnode}[5]{\treenode{#1}{#2}{\K{#3}\GB{dd}{.}\B[-8]{dddl}\B[-8]{dddr}}{#4}{#5}}
\newcommand{\dotbelownode}[5]{\treenode{#1}{#2}{\K{#3}\GB{dd}{.}}{#4}{#5}}
\newcommand{\bnode}[5]{\treenode{#1}{#2}{\K{#3}\branchbelow}{#4}{#5}}
\newcommand{\bnodeticked}[5]{\treenodeticked{#1}{#2}{\K{#3}\branchbelow}{#4}{#5}}
\newcommand{\barenode}[1]{\K{#1}}
\newcommand{\dotbelowbarenode}[1]{\K{#1}\GB{dd}{.}}
\newcommand{\dotbelowbarebnode}[1]{\K{#1}\GB{dd}{.}\B[-8]{dddl}\B[-8]{dddr}}
\newcommand{\dotbelowbaretribnode}[1]{\K{#1}\GB{dd}{.}\B[-8]{dddl}\B[-8]{ddd}\B[-8]{dddr}}

% Syntax:
%
% \tree[1]{ ... } => optional 1 adds width to branching
%
% Trees are defined like tabular tables:
%
% \tree{ & node1 & \\ node2 & & node3 }
%
% A tree node is defined by
%
% \nnode{10}{1.}{$A \land B$}{w}{Ass.}
%
% The first argument (10) controls the distance between the formula
% and the line number/world/annotation.
%
% Nodes under which the tree branches are defined by \bnode instead of
% \nnode. A bnode must be followed by an empty row:
%
%\tree[2]{
%    & \bnode{13}{1.}{$A \lor B$}{w}{Ass.} & \\
%    && \\
%    \nnode{8}{2.}{$A$}{w}{(1)} && \nnode{8}{3.}{$B$}{w}{(1)}
%}

\begin{document}

\section*{Answers to Selected Exercises}\label{ch:answers}

\subsection*{Exercise 1.1}

An operator $O$ is truth-functional if you can figure out the
truth-value of $Op$ from the truth-value of $p$. So (c) and
(g) are truth-functional; (a), (b), (d), and (e) are not
truth-functional. (f) is truth-functional if God is omniscient (and
infallible); it is also truth-functional if God doesn't exist, or if
God believes all and only false things; otherwise (f) is not
truth-functional.

\subsection*{Exercise 1.3}

(a), (b), and (e).

\subsection*{Exercise 1.4}

  The following pairs are duals: (a) and (c), (b) and (d), (e) and
  (g), (f) and (h), (i) and (k), (l) and (l), (m) and (m).

\subsection*{Exercise 1.6}

  All instances of $\Box \Box A \to \Box A$ are instances of
  \textbf{T}. So if all instances of \textbf{T} are valid, then so are
  all instances of $\Box \Box A \to \Box A$.

\subsection*{Exercise 1.7}

  I show that if $\Box A \to A$ is valid, then so is
  $A \to \Diamond A$. The other direction is similar.

  $\Box A \to A$ is equivalent to $\neg A \to \neg \Box A$, which (by
  duality) is equivalent to $\neg A \to \Diamond \neg A$. If
  $\neg A \to \Diamond \neg A$ is valid then so is
  $\neg \neg A \to \Diamond \neg \neg A$, because all instances of the
  latter are instances of the former. And
  $\neg \neg A \to \Diamond \neg \neg A$ is equivalent to
  $A \to \Diamond A$.


\subsection*{Exercise 1.12}

  \textbf{D} is valid because $A$ and $\neg A$ can never both be true
  in virtue of their form. So if $A$ is true in virtue of its form,
  then $\neg A$ is not true in virtue of its form.

  For \textbf{G}, suppose $\Diamond\Box A$. Then there is some
  scenario at which $\Box A$ is true under some interpretation of the
  non-logical vocabulary. But whether $A$ is logically true doesn't
  vary from scenario to scenario or from interpretation to
  interpretation. So $\Box A$ is true at every scenario under every
  interpretation (of the non-logical vocabulary). And then
  $\Diamond A$ is true at every scenario under every interpretation
  (as reasoned above for \textbf{D}). So $\Box\Diamond A$.


\subsection*{Exercise 1.13}

  \begin{enumerate}[(a)]
  \item All of them.
  \item \textbf{K}, \textbf{Dual1}, \textbf{Dual2}, \textbf{4},
    \textbf{5}, and \textbf{G}, but not \textbf{T} or \textbf{D}.
  \end{enumerate}

\subsection*{Exercise 2.2}

  At both worlds.

\subsection*{Exercise 2.3}

  Suppose for reductio that some instance of $\Box A \to \Diamond A$ is false at
  some world $w$ in some model. By definition 2.2, $\Box A$
  is then true at $w$ and $\Diamond A$ false. But if $\Diamond A$ is false at
  $w$ then (by definition 2.2) $A$ is false at every world
  in the model, including $w$. And then $\Box A$ isn't true at $w$ (by
  definition 2.2). Contradiction.

\subsection*{Exercise 2.4}

  \textbf{5} reduces to the tautology $\Diamond A \to \Diamond A$; \textbf{G}
  reduces to $\Box A \to \Diamond A$, whose validity we proved above.

\subsection*{Exercise 2.5}

  Suppose $A$ is valid. So $A$ is true at all worlds in all models (by
  definition 2.3). It follows that in any given model, $A$ is true
  at every world. By definition 2.2, it follows that
  $\Box A$ is true at every world in any model.

\subsection*{Exercise 2.8}

  \begin{enumerate}[(a)]

  \item Target: $p \to \Box(p \lor q)$

    \tree{
      \nnode{25}{1.}{$\neg (p \to \Box(p \lor q))$}{w}{(A)}\\
      \nnode{25}{2.}{$p$}{w}{(1)}\\
      \nnode{25}{3.}{$\neg \Box(p \lor q)$}{w}{(1)}\\
      \nnode{25}{4.}{$\neg(p \lor q)$}{v}{(3)}\\
      \nnode{25}{5.}{$\neg p$}{v}{(4)}\\
      \nnode{25}{6.}{$\neg q$}{v}{(4)}
    }

    Countermodel: $W = \{ w, v \}, V(p,w) = 1, V(p,v)=0, V(q,v)=0$.

  \medskip
  \item Target: $\Box p \lor \Box \neg p$

    \tree{
      \nnode{25}{1.}{$\neg (\Box p \lor \Box \neg p)$}{w}{(A)}\\
      \nnode{25}{2.}{$\neg \Box p$}{w}{(1)}\\
      \nnode{25}{3.}{$\neg \Box \neg p$}{w}{(1)}\\
      \nnode{25}{4.}{$\neg p$}{v}{(2)}\\
      \nnode{25}{5.}{$\neg \neg p$}{u}{(3)}\\
      \nnode{25}{6.}{$p$}{u}{(5)}
    }

    Countermodel: $W = \{ w,v,u \}, V(p,v)=0, V(p,u)=1$.

    \medskip
  \item Target: $\Diamond(p \to q) \to (\Diamond p \to \Diamond q)$

    \tree[3]{
      &\nnode{35}{1.}{$\neg(\Diamond(p \to q) \to (\Diamond p \to \Diamond q))$}{w}{(A)}&\\
      &\nnode{35}{2.}{$\Diamond(p \to q)$}{w}{(1)}&\\
      &\nnode{35}{3.}{$\neg (\Diamond p \to \Diamond q)$}{w}{(1)}&\\
      &\nnode{35}{4.}{$\Diamond p$}{w}{(3)}&\\
      &\nnode{35}{5.}{$\neg \Diamond q$}{w}{(3)}&\\
      &\nnode{35}{6.}{$p\to q$}{v}{(2)}&\\
      &\nnode{35}{7.}{$p$}{u}{(4)}&\\
      &\nnode{35}{8.}{$\neg q$}{w}{(5)}&\\
      &\nnode{35}{9.}{$\neg q$}{v}{(5)}&\\
      &\bnode{35}{10.}{$\neg q$}{u}{(5)}&\\
      &&\\
      \nnode{10}{11.}{$\neg p$}{v}{(6)} && \nnodeclosed{10}{12.}{$q$}{v}{(6)}
    }

    Countermodel: $W = \{ w,v,u \}, V(p,v)=0, V(p,u)=1, V(q,w)=0, V(q,v)=0, V(q,u)=0$.

  %\item Target: $p\to q$

  %\item $\Box \Diamond p \to p$

  \end{enumerate}

\subsection*{Exercise 2.10}

  \begin{enumerate}[(a)]

  \item Target: $p \to \Box\Diamond p$

    \tree{
      \nnode{25}{1.}{$\neg (p \to \Box \Diamond p)$}{w}{(A)}\\
      \nnode{25}{2.}{$p$}{w}{(1)}\\
      \nnode{25}{3.}{$\neg\Box\Diamond p$}{w}{(1)}\\
      \nnode{25}{4.}{$\neg\Diamond p$}{v}{(3)}\\
      \nnodeclosed{25}{5.}{$\neg p$}{w}{(4)}
    }

    The target sentence is valid.
    \medskip
  \item Target: $\Diamond\Diamond p \to \Diamond p$

    \tree{
      \nnode{25}{1.}{$\neg (\Diamond\Diamond p \to \Diamond p)$}{w}{(A)}\\
      \nnode{25}{2.}{$\Diamond\Diamond p$}{w}{(1)}\\
      \nnode{25}{3.}{$\neg\Diamond p$}{w}{(1)}\\
      \nnode{25}{4.}{$\Diamond p$}{v}{(2)}\\
      \nnode{25}{5.}{$p$}{u}{(4)}\\
      \nnodeclosed{25}{6.}{$\neg p$}{u}{(3)}
    }

    The target sentence is valid.
    \medskip


  \item Target: $\Diamond(p \land q) \to (\Diamond p \land \Diamond q)$

    \tree[3]{
      &\nnode{35}{1.}{$\neg(\Diamond(p \land q) \to (\Diamond p \land \Diamond q))$}{w}{(A)}&\\
      &\nnode{35}{2.}{$\Diamond(p \land q)$}{w}{(1)}&\\
      &\nnode{35}{3.}{$\neg(\Diamond p \land \Diamond q)$}{w}{(1)}&\\
      &\nnode{35}{4.}{$p \land q$}{v}{(2)}&\\
      &\nnode{35}{5.}{$p$}{v}{(4)}&\\
      &\bnode{35}{6.}{$q$}{v}{(4)}&\\
      &&\\
      \nnode{12}{7.}{$\neg \Diamond p$}{w}{(3)} && \nnode{12}{8.}{$\neg \Diamond q$}{w}{(3)}\\
      \nnodeclosed{12}{9.}{$\neg p$}{v}{(7)} && \nnodeclosed{12}{10.}{$\neg q$}{v}{(8)}
    }

    The target sentence is valid.
    \medskip

  \item Target: $(\Diamond p \land \Diamond q) \to \Diamond(p \land q)$

    \tree[3]{
      &\nnode{35}{1.}{$\neg ((\Diamond p \land \Diamond q) \to \Diamond(p \land q))$}{w}{(A)}&&\\
      &\nnode{35}{2.}{$\Diamond p \land \Diamond q$}{w}{(1)}&&\\
      &\nnode{35}{3.}{$\neg \Diamond(p \land q)$}{w}{(1)}&&\\
      &\nnode{35}{4.}{$\Diamond p$}{w}{(2)}&&\\
      &\nnode{35}{5.}{$\Diamond q$}{w}{(2)}&&\\
      &\nnode{35}{6.}{$p$}{v}{(4)}&&\\
      &\nnode{35}{7.}{$q$}{u}{(5)}&&\\
      &\bnode{35}{8.}{$\neg (p \land q)$}{v}{(3)}&&\\
      &&&\\
      \nnodeclosed{9}{9.}{$\neg p$}{v}{(8)}&&\nnode{15}{10.}{$\neg q$}{v}{(8)}&\\
      &&\bnode{15}{11.}{$\neg(p \land q)$}{u}{(3)}&\\
      &&&\\
      &\nnode{15}{12.}{$\neg p$}{u}{(11)}&&\nnodeclosed{9}{13.}{$\neg q$}{u}{(11)}&\\
      &\bnode{15}{14.}{$\neg(p \land q)$}{w}{(3)}&&\\
      &&&\\
      \nnode{9}{15.}{$\neg p$}{w}{(14)}&&\nnode{9}{16.}{$\neg q$}{w}{(14)}&
    }

    Countermodel: $W = \{ w,v,u \}, V(p,w)=0, V(p,v)=1, V(p,u)=0, V(q,v)=0, V(q,u)=1$.

    \medskip

  \item Target: $\Diamond(p \lor q) \leftrightarrow (\Diamond p \lor \Diamond q)$

    DIY. The sentence is valid.

    % can't draw the tree!

    % \tree[6]{
    %   &&&\bnode{30}{1.}{$\Diamond(p \lor q) \leftrightarrow (\Diamond p \lor \Diamond q)$}{w}{(A)}&&\\
    %   &&&&&\\
    %   &&\nnode{20}{2.}{$\Diamond(p \lor q)$}{w}{(1)} && \nnode{20}{4.}{$\neg\Diamond(p \lor q)$}{w}{(1)}&\\
    %   &&\nnode{20}{3.}{$\neg(\Diamond p \lor \Diamond q)$}{w}{(1)}&&\bnode{20}{5.}{$\Diamond p \lor \Diamond q$}{w}{(1)}&\\
    %   &&\nnode{20}{6.}{$p \lor q)$}{v}{(2)}&&&\\
    %   &&\nnode{20}{7.}{$\neg \Diamond p)$}{w}{(4)}& \nnode{10}{13.}{$\Diamond p$}{w}{(5)}&& \nnode{10}{14.}{$\Diamond q$}{w}{(5)}\\
    %   &&\bnode{20}{8.}{$\neg \Diamond q)$}{w}{(4)}& \nnode{10}{15.}{$p$}{v}{(13)}&& \nnode{10}{19.}{$q$}{v}{(14)}\\
    %   &&& \nnode{10}{15.}{$\neg (p\lor q)$}{v}{(4)}&& \nnode{10}{20.}{$\neg(p \lor q)$}{v}{(4)}\\
    %   &&& \nnode{10}{15.}{$\neg (p\lor q)$}{v}{(4)}&& \nnode{10}{20.}{$\neg(p \lor q)$}{v}{(4)}\\
    % }

    \medskip

  \item Target: $\Box\Diamond p \to \Diamond\Box p$

    \tree[3]{
      &\nnode{25}{1.}{$\neg (\Box\Diamond p \to \Diamond\Box p)$}{w}{(A)}&&\\
      &\nnode{25}{2.}{$\Box\Diamond p$}{w}{(1)}&&\\
      &\nnode{25}{3.}{$\neg \Diamond\Box p$}{w}{(1)}&&\\
      &\nnode{25}{4.}{$\Diamond p$}{w}{(2)}&&\\
      &\nnode{25}{5.}{$p$}{v}{(4)}&&\\
      &\nnode{25}{6.}{$\neg \Box p$}{w}{(3)}&&\\
      &\nnode{25}{7.}{$\neg p$}{u}{(6)}&&\\
      &\nnode{25}{8.}{$\Diamond p$}{v}{(2)}&&\\
      &\nnode{25}{9.}{$p$}{s}{(8)}&&\\
      &\nnode{25}{10.}{$\neg \Box p$}{v}{(3)}&&\\
      &\nnode{25}{11.}{$\neg p$}{t}{(10)}&&\\
      &\nnode{25}{}{\vdots}{}{}&&
    }

    This tree grows forever. Manually inspecting the countermodel we
    could read off from the tree at step 7 shows that it is a genuine
    countermodel.

    Countermodel: $W = \{ w,v,u \}, V(p,v)=1, V(p,u)=0$.

  \end{enumerate}

\subsection*{Exercise 2.11}

% ------------------------------------------------------------------------
% file `ml-exercise-24-solution-body.tex'
% ------------------------------------------------------------------------
%
  We could show that the conditional $(A_1\land\ldots\land A_n) \to B$ is
  valid. More simply, we can start a tree with individual nodes for the premises
  and the negated conclusion. If the tree closes, $A_1,\ldots,A_n$ entail $B$.

\subsection*{Exercise 3.1}

  $v$ has access to no world. So any sentence $A$ is true at \emph{all} (zero)
  worlds accessible from $v$.

  If this seems strange, remember that $\Box A$ is equivalent to
  $\neg \Diamond \neg A$. And $\Diamond \neg A$ means that there's an
  accessible world where $\neg A$ is true. If there are no accessible
  worlds, then this is false. So $\neg \Diamond \neg A$ is true.


\subsection*{Exercise 3.3}

  There are infinitely many correct answers for each world. For
  example: $w_1: \Diamond\Box p$, $w_2: \neg p \land \neg q$,
  $w_3: \Box p$, $w_4: \Box q$.

\subsection*{Exercise 3.5}

  For example: $W = \{ w \}$, not $w R w$, $V(p,w) = 1$.

\subsection*{Exercise 3.6}

  For example: $\Box(p \lor \neg p) \to (p \lor \neg p)$.

\subsection*{Exercise 3.8}

  Reflexive yes,  serial yes, transitive yes,  euclidean no, symmetric
  no, universal no.

\subsection*{Exercise 3.10}

% ------------------------------------------------------------------------
% file `ml-exercise-34-solution-body.tex'
%
%     solution of type `exercise' with id `34'
%
% generated by the `solution' environment of the
%   `xsim' package v0.16a (2020/01/16)
% from source `ml' on 2020/09/26 on line 859
% ------------------------------------------------------------------------
  It's true that if $R$ is symmetric and transitive then $wRv$ implies
  $vRw$ which implies $wRw$. But this only shows that every world $w$
  \emph{that can see some world $v$} can see itself. What's true is
  that symmetry, transitivity, \emph{and seriality} together imply
  reflexivity.

\subsection*{Exercise 3.11}

% ------------------------------------------------------------------------
% file `ml-exercise-35-solution-body.tex'
%
%     solution of type `exercise' with id `35'
%
% generated by the `solution' environment of the
%   `xsim' package v0.16a (2020/01/16)
% from source `ml' on 2020/09/26 on line 1060
% ------------------------------------------------------------------------
  \begin{enumerate}[(a)]
  \item Every world has access only to itself.
  \item No world has access to any world.
\end{enumerate}

\subsection*{Exercise 4.3}

% ------------------------------------------------------------------------
% file `ml-exercise-39-solution-body.tex'
%
%     solution of type `exercise' with id `39'
%
% generated by the `solution' environment of the
%   `xsim' package v0.16a (2020/01/16)
% from source `ml' on 2020/09/26 on line 446
% ------------------------------------------------------------------------
  We have to show that if there some K-tree for a given target sentence closes,
  then there is no fully expanded K-tree for the sentence that remains open.
  So assume some K-tree for the target sentence closes. By the Completeness
  Theorem, the target sentence is K-valid. If there were a fully expanded but
  open K-tree for the same target sentence, then by the Completeness Lemma the
  negation of that sentence would be true at world $w$ in the model induced by
  some open branch on that tree. This contradicts the fact that the target
  sentence is K-valid.

  \subsection*{Exercise 4.4}

  % ------------------------------------------------------------------------
% file `ml-exercise-40-solution-body.tex'
%
%     solution of type `exercise' with id `40'
%
% generated by the `solution' environment of the
%   `xsim' package v0.16a (2020/01/16)
% from source `ml' on 2020/09/26 on line 617
% ------------------------------------------------------------------------
  \begin{alignat*}{2}
    1.\quad& \Box \neg p \to \neg p &\quad& \text{(T)}\\
    2.\quad& (\Box \neg p \to \neg p) \to (p \to \neg \Box \neg p) &\quad& \text{(Taut)}\\
    3.\quad& p \to \neg \Box \neg p &\quad& \text{(1, 2, MP)}\\
    4.\quad& \Box p \to p &\quad& \text{(T)}\\
    5.\quad& (\Box p \to p) \to ((p \to \neg \Box \neg p) \to (\Box p \to \neg \Box \neg p)) &\quad& \text{(Taut)}\\
    6.\quad& (p \to \neg \Box \neg p) \to (\Box p \to \neg \Box \neg p) &\quad& \text{(4, 5, MP)}\\
    7.\quad& \Box p \to \neg \Box \neg p &\quad& \text{(3, 6, MP)}
  \end{alignat*}

  \subsection*{Exercise 4.5}

% ------------------------------------------------------------------------
% file `ml-exercise-41-solution-body.tex'
%
%     solution of type `exercise' with id `41'
%
% generated by the `solution' environment of the
%   `xsim' package v0.16a (2020/01/16)
% from source `ml' on 2020/09/26 on line 641
% ------------------------------------------------------------------------
  \begin{alignat*}{2}
    1.\quad& A \to \Diamond A &\quad& \text{(\textbf{T}, Prop.\ Logic)}\\
    2.\quad& \Diamond A \to \Box \Diamond A &\quad& \text{(\textbf{5})}\\
    3.\quad& \Box\Diamond A \to \Diamond\Box A &\quad& \text{(\textbf{M})}\\
    4.\quad& \Diamond\Box A \to \Box A &\quad& \text{(\textbf{5}, Prop.\ Logic)}\\
    5.\quad& A \to \Box A &\quad& \text{(1, 2, 3, 4, Prop.\ Logic)}
  \end{alignat*}

  \subsection*{Exercise 4.7}

% ------------------------------------------------------------------------
% file `ml-exercise-43-solution-body.tex'
%
%     solution of type `exercise' with id `43'
%
% generated by the `solution' environment of the
%   `xsim' package v0.16a (2020/01/16)
% from source `ml' on 2020/09/26 on line 759
% ------------------------------------------------------------------------
  We need to show that everything that's derivable in the axiomatic calculus for
  S4 is true at every world in every transitive and reflexive Kripke model. From
  the soundness proof for K, we know that all instances of
  \textbf{A1}--\textbf{A3} are true at every world in every Kripke model. From
  observation 3.3, we know that all instances of \textbf{T} are
  true at every world in every reflexive Kripke model. From observation
  3.4, we know that all instances of \textbf{4} are true at every
  world in every transitive Kripke model. So all axioms in the S4-calculus are
  valid in the class of transitive and reflexive Kripke frames. Since
  \textbf{MP} and \textbf{Nec} preserve validity in any class of Kripke frames,
  it follows that everything that's derivable in the S4-calculus is valid in the
  class of transitive and reflexive frames.

  \subsection*{Exercise 4.8}

% ------------------------------------------------------------------------
% file `ml-exercise-44-solution-body.tex'
%
%     solution of type `exercise' with id `44'
%
% generated by the `solution' environment of the
%   `xsim' package v0.16a (2020/01/16)
% from source `ml' on 2020/09/26 on line 1103
% ------------------------------------------------------------------------
  We have to show that all S4-valid sentences are provable in the axiomatic
  calculus for S4, which extends the calculus for T by the axiom (schema)
  $\Box A \to \Box\Box A$. The argument is by contraposition: We suppose that
  some sentence is not S4-provable and show that it is not S4-valid.

  Suppose $A$ is not S4-provable. Then $\{ \neg A \}$ is S4-consistent. It
  follows by Lindenbaum's Lemma that $\{ \neg A \}$ is included in some maximal
  S4-consistent set $S$. By definition of canonical models, that set is a world
  in the canonical model $M_{S4}$ for S4.  Since $\neg A$ is in $S$, it follows
  from the Canonical Model Lemma that $M_{S4},S \models \neg A$. So
  $M_{S4},S \not\models A$.

  It remains to show that $M_{S4}$ is reflexive and transitive.

  By definition, a world $v$ in a canonical model is accessible from $w$ iff
  whenever $\Box A \in w$ then $A \in v$. Since the worlds in $M_{S4}$ are
  maximal S4-consistent sets of sentences, and every such set contains every
  instance of the \textbf{T} schema $\Box A \to A$, there is no world in
  $M_{S4}$ that contains $\Box A$ but not $A$. So every world in $M_{S4}$ has
  access to itself. So $M_{S4}$ is reflexive.

  For transitivity, suppose for some worlds $w,v,u$ in $M_{S4}$ we have $wRv$
  and $vRu$. We need to show that $wRu$. Given how $R$ is defined in $M_{S4}$,
  we have to show that $u$ contains all sentences $A$ for which $w$ contains
  $\Box A$. So let $A$ be an arbitrary sentence for which $w$ contains $\Box
  A$. Since every world in $M_{S4}$ contains every instance of
  $\Box A \to \Box\Box A$, we know that $w$ also contains $\Box\Box A$. From
  $wRv$, we can infer that $v$ contains $\Box A$. And from $vRu$, we can infer
  that $u$ contains $A$.

  \subsection*{Exercise 4.9}

% ------------------------------------------------------------------------
% file `ml-exercise-45-solution-body.tex'
%
%     solution of type `exercise' with id `45'
%
% generated by the `solution' environment of the
%   `xsim' package v0.16a (2020/01/16)
% from source `ml' on 2020/09/26 on line 1402
% ------------------------------------------------------------------------
  The argument closely follows the argument I gave on
  p.87. The \textbf{5} schema would give us
  $\neg \Box \neg A \to \Box \neg \Box \neg A$. In the minimal logic
  K, the consequent $\Box \neg \Box \neg A$ entails
  $\Box(\Box \neg A \to \neg A)$, which entails $\Box \neg A$ by the
  \textbf{GL} schema and \textbf{MP}. So if the \textbf{5} schema were
  valid in the logic of provability, then the schema
  $\neg \Box \neg A \to \Box \neg A$ would also be valid. And that
  schema is only valid if everything whatsoever is provable, even
  contradictions.

\subsection*{Exercise 5.2}

% ------------------------------------------------------------------------
% file `ml-exercise-47-solution-body.tex'
%
%     solution of type `exercise' with id `47'
%
% generated by the `solution' environment of the
%   `xsim' package v0.16a (2020/01/16)
% from source `ml' on 2020/09/26 on line 269
% ------------------------------------------------------------------------
  \begin{enumerate}[(a)]
  \item $\Kn_a(r \lor s)$\\
    $r$: It is raining; $s$: It is snowing\\[-2mm]
  \item $\Kn_a r \lor \Kn_a s$\\
    $r$: It is raining; $s$: It is snowing\\[-2mm]
  \item $\Kn_b r \lor \Kn_b \neg r$\\
    $r$: It is raining\\[-2mm]
  \item $\Kn_c \neg \Kn_c r$\\
    $r$: It is raining\\[-2mm]
  \item $\Kn_a(\Kn_b r \lor \Kn_b \neg r)$\\
    $r$: It is raining
  \end{enumerate}

\subsection*{Exercise 5.3}

% ------------------------------------------------------------------------
% file `ml-exercise-48-solution-body.tex'
%
%     solution of type `exercise' with id `48'
%
% generated by the `solution' environment of the
%   `xsim' package v0.16a (2020/01/16)
% from source `ml' on 2020/09/26 on line 308
% ------------------------------------------------------------------------
  One possible answer: There are many propositions that I don't know, and that
  don't logically follow from things I know. And for many of these propositions
  I \emph{know} that that they don't follow from things I know. (For example, I
  know that it doesn't follow from anything I know that Edinburgh is in Italy.)
  Let $p$ be some such proposition. Since I know that $\neg \Kn^+ p$ is true, I
  can easily figure out that $\Kn^+ p \to p$ is true as well (by the truth-table
  for the arrow). So $\Kn^*(\Kn^+p \to p)$. If
  $\Kn^*(\Kn^+ A \to A) \to \Kn^+ A$ were valid, it would follow that $p$ does
  after all follow from what I know!


  \subsection*{Exercise 5.4}

  % ------------------------------------------------------------------------
% file `ml-exercise-49-solution-body.tex'
%
%     solution of type `exercise' with id `49'
%
% generated by the `solution' environment of the
%   `xsim' package v0.16a (2020/01/16)
% from source `ml' on 2020/09/26 on line 598
% ------------------------------------------------------------------------
  see
  \href{https://plato.stanford.edu/entries/dynamic-epistemic/appendix-B-solutions.html}{https://plato.stanford.edu/entries/dynamic-epistemic/appendix-B-solutions.html}


  \subsection*{Exercise 5.7}

% ------------------------------------------------------------------------
% file `ml-exercise-52-solution-body.tex'
%
%     solution of type `exercise' with id `52'
%
% generated by the `solution' environment of the
%   `xsim' package v0.16a (2020/01/16)
% from source `ml' on 2020/09/26 on line 950
% ------------------------------------------------------------------------
  You can find the trees for (a)--(d) on
  \href{https://www.wolfgangschwarz.net/trees/}{https://www.wolfgangschwarz.net/trees/}
  if you write $\Kn$ as a box and $\Mi$ as a diamond.

  \subsection*{Exercise 5.10}

% ------------------------------------------------------------------------
% file `ml-exercise-55-solution-body.tex'
%
%     solution of type `exercise' with id `55'
%
% generated by the `solution' environment of the
%   `xsim' package v0.16a (2020/01/16)
% from source `ml' on 2020/09/26 on line 1027
% ------------------------------------------------------------------------
  (a) and (b) are valid, (c) and (d) are invalid. One way of testing
  this is with the method of trees; but you'll need to keep track of
  both accessibility relations. To illustrate, here is a tree proof
  for (a).

  \bigskip
  \begin{center}
  \tree{
    \nnode{25}{1.}{$\neg(\Mi_1 \Kn_2 A \to \Mi_1 A)$}{w}{(Ass.)}\\
    \nnode{25}{2.}{$\Mi_1 \Kn_2 A$}{w}{(1)}\\
    \nnode{25}{3.}{$\neg \Mi_1 A$}{w}{(1)}\\
    \nnode{25}{4.}{$wR_1 v$}{}{(2)}\\
    \nnode{25}{5.}{$\Kn_2 A$}{v}{(2)}\\
    \nnode{25}{6.}{$\neg A$}{v}{(3,4)}\\
    \nnode{25}{7.}{$vR_2 v$}{}{\;\;(Refl.)}\\
    \nnodeclosed{25}{8.}{$A$}{v}{(5,7)}
  }
  \end{center}

  The tree for (c) doesn't close:

  \bigskip
  \begin{center}
  \tree{
    \nnode{25}{1.}{$\neg(\Mi_1 \Kn_2 A \to \Mi_2\Kn_1 A)$}{w}{(Ass.)}\\
    \nnode{25}{2.}{$\Mi_1 \Kn_2 A$}{w}{(1)}\\
    \nnode{25}{3.}{$\neg \Mi_2\Kn_1 A$}{w}{(1)}\\
    \nnode{25}{4.}{$wR_1 v$}{}{(2)}\\
    \nnode{25}{5.}{$\Kn_2 A$}{v}{(2)}\\
    \nnode{25}{6.}{$vR_2v$}{}{\;\;(Refl.)}\\
    \nnode{25}{7.}{$A$}{v}{(5,6)}\\
    \nnode{25}{8.}{$wR_2 w$}{}{\;\;(Refl.)}\\
    \nnode{25}{9.}{$\neg\Kn_1 A$}{w}{(3,8)}\\
    \nnode{25}{10.}{$wR_1u$}{}{(9)}\\
    \nnode{25}{11.}{$\neg A$}{u}{(9)}
  }
  \end{center}
  %
  We could add a few more applications of Reflexivity, but the tree
  would remain open. It also gives us a countermodel: let $W$ =
  $\{ w,v,u \}$; $w$ has 1-access to $v$ and $u$; each world has 1-
  and 2-access to itself; $V(p,v) = 1$ and $V(p,u) = 0$. In this
  model, at world $w$, $\Mi_1\Kn_2 p$ is true while $\Mi_2\Kn_1 p$ is
  false.

  Cases (b) and (d) are similar.

  \subsection*{Exercise 5.11}

 % ------------------------------------------------------------------------
% file `ml-exercise-56-solution-body.tex'
%
%     solution of type `exercise' with id `56'
%
% generated by the `solution' environment of the
%   `xsim' package v0.16a (2020/01/16)
% from source `ml' on 2020/09/26 on line 1098
% ------------------------------------------------------------------------
  The \textbf{5}-schema for $\EKn_G$ states that
  $\neg \EKn_G \neg A \to \EKn_G \neg \EKn_G \neg A$. To show that
  this is invalid, we need to find a case where some instance of
  $\neg \EKn_G \neg A$ is true while $\EKn_G \neg \EKn_G \neg A$ is
  false. Let's take the simplest instance, with $A=p$, and let's
  assume the relevant group has two agents. Consider a world $w$ at
  which $\Kn_1 \neg p$ and $\neg \Kn_2 \neg p$ are true. By the
  assumption that \textbf{5} is valid for $\Kn_i$,
  $\Kn_2\neg\Kn_2\neg p$ is also true at $w$. But
  $\Kn_1\neg \Kn_2\neg p$ can be false (at $w$). If it is, then
  $\neg \EKn_G \neg p$ is true at $w$ while
  $\EKn_G \neg \EKn_G \neg p$ is false.

  \subsection*{Exercise 5.12}

% ------------------------------------------------------------------------
% file `ml-exercise-57-solution-body.tex'
%
%     solution of type `exercise' with id `57'
%
% generated by the `solution' environment of the
%   `xsim' package v0.16a (2020/01/16)
% from source `ml' on 2020/09/26 on line 1217
% ------------------------------------------------------------------------
  No, a transitive, serial, and euclidean relation is not always
  symmetric. (Counterexample: wRv, vRv.) This means that \textbf{B} (which
  corresponds to symmetry) is not valid in KD45.


  \subsection*{Exercise 5.13}

  % ------------------------------------------------------------------------
% file `ml-exercise-58-solution-body.tex'
%
%     solution of type `exercise' with id `58'
%
% generated by the `solution' environment of the
%   `xsim' package v0.16a (2020/01/16)
% from source `ml' on 2020/09/26 on line 1227
% ------------------------------------------------------------------------
  You can e.g. do a tree proof, using $\Bel$ as the box.

  \subsection*{Exercise 5.14}

% ------------------------------------------------------------------------
% file `ml-exercise-59-solution-body.tex'
%
%     solution of type `exercise' with id `59'
%
% generated by the `solution' environment of the
%   `xsim' package v0.16a (2020/01/16)
% from source `ml' on 2020/09/26 on line 1329
% ------------------------------------------------------------------------
  By \textbf{PI}, $\Bel A \to \Kn\Bel A$ is valid. By \textbf{KB}, so
  is $\Kn\Bel A \to \Bel\Bel A$. By propositional logic, it follows
  that $\Bel A \to \Bel\Bel A$ is valid, which is principle \textbf{4}
  for belief.

  By \textbf{NI}, $\neg \Bel \neg A \to \Kn\neg \Bel \neg A$ is valid. By \textbf{KB}, so
  is $\Kn\neg \Bel \neg A \to \Bel\neg \Bel \neg A$. By propositional logic, it follows
  that $\neg \Bel \neg A \to \Bel\neg \Bel \neg A$ is valid, which is principle \textbf{5}
  for belief.

  \subsection*{Exercise 6.1}

 % ------------------------------------------------------------------------
% file `ml-exercise-63-solution-body.tex'
%
%     solution of type `exercise' with id `63'
%
% generated by the `solution' environment of the
%   `xsim' package v0.16a (2020/01/16)
% from source `ml' on 2020/09/26 on line 143
% ------------------------------------------------------------------------
  For example: $\Ob \neg$, or $\neg \Pe$.

  \subsection*{Exercise 6.3}
  % ------------------------------------------------------------------------
% file `ml-exercise-65-solution-body.tex'
%
%     solution of type `exercise' with id `65'
%
% generated by the `solution' environment of the
%   `xsim' package v0.16a (2020/01/16)
% from source `ml' on 2020/09/26 on line 326
% ------------------------------------------------------------------------
  (a): $\Box(\OK \to (\Box(\OK \to A) \to A))$.

  \subsection*{Exercise 6.4}

  % ------------------------------------------------------------------------
% file `ml-exercise-66-solution-body.tex'
%
%     solution of type `exercise' with id `66'
%
% generated by the `solution' environment of the
%   `xsim' package v0.16a (2020/01/16)
% from source `ml' on 2020/09/26 on line 339
% ------------------------------------------------------------------------
  $\Pe A$ could be defined as $\neg\Box(\OK->\neg A)$, or more simply (and
  equivalently) as $\Diamond (\OK \land A)$.

  \subsection*{Exercise 6.5}

  % ------------------------------------------------------------------------
% file `ml-exercise-67-solution-body.tex'
%
%     solution of type `exercise' with id `67'
%
% generated by the `solution' environment of the
%   `xsim' package v0.16a (2020/01/16)
% from source `ml' on 2020/09/26 on line 409
% ------------------------------------------------------------------------
  On the absolutist conception, $wRv$ means that $v$ is ideal. So transitivity
  (if $wRv$ and $vRu$ then $wRu$) and euclidity (if $wRv$ and $wRu$ then $vRu$)
  both state that if $v$ is ideal and $u$ is ideal then $u$ is ideal.

  \subsection*{Exercise 6.6}

  % ------------------------------------------------------------------------
% file `ml-exercise-68-solution-body.tex'
%
%     solution of type `exercise' with id `68'
%
% generated by the `solution' environment of the
%   `xsim' package v0.16a (2020/01/16)
% from source `ml' on 2020/09/26 on line 429
% ------------------------------------------------------------------------
  Consider the example from the text, where $w$ is the actual world (in the UK)
  and $u$ is a $w$-accessible world at which everyone drives on the left
  although the law says that one must drive on the right. A typical world
  accessible from $u$ will be a world at which people drive on the right. This
  world will not be accessible from $w$. So we have a counterexample to
  transitivity. We also have a counterexample to euclidity because we have $wRu$
  and $wRu$ but not $uRu$. (Euclidity entails shift reflexivity.)

  \subsection*{Exercise 6.8}

  % ------------------------------------------------------------------------
% file `ml-exercise-70-solution-body.tex'
%
%     solution of type `exercise' with id `70'
%
% generated by the `solution' environment of the
%   `xsim' package v0.16a (2020/01/16)
% from source `ml' on 2020/09/26 on line 493
% ------------------------------------------------------------------------
  In the described situation, it ought to be the case that Amy is either
  obligated to help Betty or obligated to help Carla, but Amy is neither
  obligated to help Betty nor to help Carla. So if $p$ translates `Amy helps
  Betty' and $q$ `Amy helps Carla', we seem to have $\Ob(\Ob p \lor \Ob q)$ and
  $\neg \Ob p$ and $\neg \Ob q$. But these assumptions are inconsistent in K5.
  You can draw a K5-tree (using the Euclidity rules) starting with
  $\Ob(\Ob p \lor \Ob q)$ and $\neg \Ob p$ and $\neg \Ob q$ on which all
  branches close. This shows that there is no world in any euclidean model at
  which the three assumptions are true.

  \subsection*{Exercise 6.9}

  % ------------------------------------------------------------------------
% file `ml-exercise-71-solution-body.tex'
%
%     solution of type `exercise' with id `71'
%
% generated by the `solution' environment of the
%   `xsim' package v0.16a (2020/01/16)
% from source `ml' on 2020/09/26 on line 526
% ------------------------------------------------------------------------
  \begin{enumerate}[(a)]
  \item I argue by contraposition. Suppose $\Ob\Ob A \to \Ob A$ is invalid
    on a frame. This means that at some world $w$ in some model $M$ based on the
    frame, (some instance of) $\Ob\Ob A$ is true while $\Ob A$ is false. It
    follows that there is a world accessible from $w$ at which $A$ is false and
    $\Ob A$ true. So $\Ob A \to A$ is false at $v$. So $\Ob (\Ob A \to A)$ is
    false at $w$. (You could also give a tree proof with the K-rules showing
    that \textbf{U} entails \textbf{C4}.)

    \item It is not enough to give a model in which some instance of \textbf{C4}
    is true at some world while the corresponding instance of \textbf{U} is
    false. For a counterexample, you need to give a \emph{frame} on which every
    instance of \textbf{C4} is valid but not every instance of \textbf{U}. Here
    is one such frame: $W = \{w,v\}$, $wRw$, $wRv$, and $vRw$.
  \end{enumerate}

  \subsection*{Exercise 6.10}

% ------------------------------------------------------------------------
% file `ml-exercise-72-solution-body.tex'
%
%     solution of type `exercise' with id `72'
%
% generated by the `solution' environment of the
%   `xsim' package v0.16a (2020/01/16)
% from source `ml' on 2020/09/26 on line 560
% ------------------------------------------------------------------------
  Use
  \href{https://www.wolfgangschwarz.net/trees/}{https://www.wolfgangschwarz.net/trees/}.
  (Write $\Ob$ as a box and $\Pe$ as a diamond, and make the accessibility
  relation serial.)
  
  \subsection*{Exercise 6.14}

  % ------------------------------------------------------------------------
% file `ml-exercise-76-solution-body.tex'
%
%     solution of type `exercise' with id `76'
%
% generated by the `solution' environment of the
%   `xsim' package v0.16a (2020/01/16)
% from source `ml' on 2020/09/26 on line 1187
% ------------------------------------------------------------------------
  Simply replace `all' in the semantics for $\Ob(B/A)$ with `some'.

  \subsection*{Exercise 6.15}

  % ------------------------------------------------------------------------
% file `ml-exercise-77-solution-body.tex'
%
%     solution of type `exercise' with id `77'
%
% generated by the `solution' environment of the
%   `xsim' package v0.16a (2020/01/16)
% from source `ml' on 2020/09/26 on line 1423
% ------------------------------------------------------------------------
  Ross's Paradox: `Alice must be in the kitchen or in the bathroom'
  seems to imply that Alice might be in the kitchen and that she might
  be in the bathroom.

  The Paradox of Free Choice: `Alice might be in the kitchen or in the
  bathroom' seems to imply that Alice might be in the kitchen and that
  she might be in the bathroom.

  \subsection*{Exercise 6.16}

  % ------------------------------------------------------------------------
% file `ml-exercise-78-solution-body.tex'
%
%     solution of type `exercise' with id `78'
%
% generated by the `solution' environment of the
%   `xsim' package v0.16a (2020/01/16)
% from source `ml' on 2020/09/26 on line 1529
% ------------------------------------------------------------------------
  Whenever $X \in N(w)$ then all sets that have $X$ as a subset are in
  $N(w)$.

  \subsection*{Exercise 6.17}

  % ------------------------------------------------------------------------
% file `ml-exercise-79-solution-body.tex'
%
%     solution of type `exercise' with id `79'
%
% generated by the `solution' environment of the
%   `xsim' package v0.16a (2020/01/16)
% from source `ml' on 2020/09/26 on line 1588
% ------------------------------------------------------------------------
  For every world $w$, every member of $N(w)$ contains $w$.

  \subsection*{Exercise 6.18}

  % ------------------------------------------------------------------------
% file `ml-exercise-80-solution-body.tex'
%
%     solution of type `exercise' with id `80'
%
% generated by the `solution' environment of the
%   `xsim' package v0.16a (2020/01/16)
% from source `ml' on 2020/09/26 on line 1628
% ------------------------------------------------------------------------
  In Kripke semantics, $\Box p$ and $\Box q$ together entail $\Box(p \land
  q)$. But if the probability of $p$ is above the threshold and the probability
  of $q$ is above the threshold, it does not follow that the probability of
  $p\land q$ is above the threshold. For example, we could have Pr$(p)=0.95$,
  Pr$(q)=0.94$, and Pr$(p \land q) = 0.95 \times 0.94 = 0.893$.

  \subsection*{Exercise 7.1}

  % ------------------------------------------------------------------------
% file `ml-exercise-81-solution-body.tex'
%
%     solution of type `exercise' with id `81'
%
% generated by the `solution' environment of the
%   `xsim' package v0.16a (2020/01/16)
% from source `ml' on 2020/09/26 on line 88
% ------------------------------------------------------------------------
  \begin{enumerate}[(a)]
  \item $\tH \neg p$\\
    $p$: It is warm\\[-2mm]
  \item $\tF p$\\
    $p$: There is a sea battle\\[-2mm]
  \item $\tF\neg \tP p$ (or, perhaps, $\neg \tF \tP p$)\\
    $p$: There is a sea battle\\[-2mm]
  \item  $\tF(p \lor \tP q)$\\
    $p$: It is warm\\[-2mm]
  \item $\neg \tP p \to \neg \tF q$\\
    $p$: You study, $q$: you pass the exam\\[-2mm]
  \item $\tP(p \land q)$\\
    $p$: I am having tea, $q$: the door bell rings
  \end{enumerate}

  \subsection*{Exercise 7.2}

  % ------------------------------------------------------------------------
% file `ml-exercise-82-solution-body.tex'
%
%     solution of type `exercise' with id `82'
%
% generated by the `solution' environment of the
%   `xsim' package v0.16a (2020/01/16)
% from source `ml' on 2020/09/26 on line 313
% ------------------------------------------------------------------------
  (a), (c), (f), (g), and (h) are true, (b), (d), and (e) are false.

  \subsection*{Exercise 7.5}

  % ------------------------------------------------------------------------
% file `ml-exercise-85-solution-body.tex'
%
%     solution of type `exercise' with id `85'
%
% generated by the `solution' environment of the
%   `xsim' package v0.16a (2020/01/16)
% from source `ml' on 2020/09/26 on line 560
% ------------------------------------------------------------------------
  Since $\tG$ is the backward-looking box, the \textbf{4} schema for $\tG$
  corresponds to transitivity of the `later than' relation $\succ$: if
  $t \succ s$ and $s\succ r$ then $t \succ r$. We can rewrite `$t \succ s$' as
  `$s\prec t$'. Similarly, `$s\succ r$' is equivalent to `$r \prec s$' and
  `$t \succ r$' to `$r \prec t$'. So transitivity of $\succ$ requires that if
  $s \prec t$ and $r\prec s$, then $r\prec t$. This evidently follows from
  transitivity of $\prec$.

  \subsection*{Exercise 7.7}

  % ------------------------------------------------------------------------
% file `ml-exercise-87-solution-body.tex'
%
%     solution of type `exercise' with id `87'
%
% generated by the `solution' environment of the
%   `xsim' package v0.16a (2020/01/16)
% from source `ml' on 2020/09/26 on line 673
% ------------------------------------------------------------------------
  If time is transitive and circular, then it is neither asymmetric nor
  irreflexive.

  \subsection*{Exercise 7.10}

 % ------------------------------------------------------------------------
% file `ml-exercise-90-solution-body.tex'
%
%     solution of type `exercise' with id `90'
%
% generated by the `solution' environment of the
%   `xsim' package v0.16a (2020/01/16)
% from source `ml' on 2020/09/26 on line 870
% ------------------------------------------------------------------------
  \begin{enumerate}[(a)]
  \item For example, $\tG A \to \tF A$.
  \item For example, $\tH A \to \tP A$.
  \item No schema corresponds to the class of frames with a last time. If we
    also assume transitivity and weak connectedness (see page
    \pageref{weakconnectedness}), then
    $\tG(A \land \neg A) \lor \tF\tG(A \land \neg A)$ works. % xxx why does
                                                             % Garanko say
                                                             % "assuming
                                                             % irreflexivity"?
  \item No schema corresponds to the class of frames with a first time. If we
    also assume transitivity and weak connectedness, then
    $\tH(A \land \neg A) \lor \tP\tH(A \land \neg A)$ works.
  \end{enumerate}

  \subsection*{Exercise 7.11}

  % ------------------------------------------------------------------------
% file `ml-exercise-91-solution-body.tex'
%
%     solution of type `exercise' with id `91'
%
% generated by the `solution' environment of the
%   `xsim' package v0.16a (2020/01/16)
% from source `ml' on 2020/09/26 on line 891
% ------------------------------------------------------------------------
  We have to show that $\tF A \to \tF\tF A$ is valid on all and only the dense
  frames. A frame is dense if whenever $t<s$ then there is a point $r$ such that
  $t<r$ and $r<s$.

  So assume a frame is dense and suppose for reductio that some instance of
  $\tF A \to \tF\tF A$ is false at some point $t$ in some model $M$ based on
  that frame. Then $\tF A$ is true at $t$ and $\tF\tF A$ is false. Since $\tF A$
  is true at $t$, it follows by definition 7.2 that $A$
  is true at some point $s$ such that $t<s$. By density, there is a point $r$
  such that $t<r<s$. But since $A$ is true at $s$, $\tF A$ is true at $r$, and
  so $\tF\tF A$ is true at $t$; contradiction.

  In the other direction, we have to show that if a frame isn't dense then some
  instance of $\tF A \to \tF\tF A$ is false at some point $t$ in some model $M$
  based on that frame. We take the simplest instance $\tF p \to \tF\tF p$. If a
  frame isn't dense then there are points $t,s$ such that $t<s$ and no point
  lies in between $t$ and $s$. Let $V$ be an interpretation function that makes
  $p$ true at $s$ and false everywhere else. Then $\tF p$ is true at $t$ but
  $\tF\tF p$ is false. So $\tF p \to \tF\tF p$ is false at $t$.

  \subsection*{Exercise 7.12}

  % ------------------------------------------------------------------------
% file `ml-exercise-92-solution-body.tex'
%
%     solution of type `exercise' with id `92'
%
% generated by the `solution' environment of the
%   `xsim' package v0.16a (2020/01/16)
% from source `ml' on 2020/09/26 on line 924
% ------------------------------------------------------------------------
  Without assumptions about the flow of time there is no way to express in
  $\L_T$ that $p$ is true at all times (or at some time). In linear flows,
  $p \land \tH p \land tG p$ does the job.

  \subsection*{Exercise 7.14}

  % ------------------------------------------------------------------------
% file `ml-exercise-94-solution-body.tex'
%
%     solution of type `exercise' with id `94'
%
% generated by the `solution' environment of the
%   `xsim' package v0.16a (2020/01/16)
% from source `ml' on 2020/09/26 on line 1183
% ------------------------------------------------------------------------
  (a)--(d) are valid, (e) is invalid.

  To show that a schema is valid, assume for reductio that there is some time
  $t$ on some history $H$ in some model $M$ at which the schema is false. Then
  (repeatedly) use definition \ref{def:ockhamistsemantics} to derive a
  contradiction.

  For (e), consider a model with three times $t,s,r$ such that $s\prec t$,
  $r\prec t$, and neither $s \prec r$ nor $r\prec s$. Let $q$ be true at $s$ and
  false at the other two times. $\tP q \to \Box \tP \Diamond q$ is false at $t$
  on the history $(s,t)$.


  \subsection*{Exercise 7.16}

  % ------------------------------------------------------------------------
% file `ml-exercise-96-solution-body.tex'
%
%     solution of type `exercise' with id `96'
%
% generated by the `solution' environment of the
%   `xsim' package v0.16a (2020/01/16)
% from source `ml' on 2020/09/26 on line 1272
% ------------------------------------------------------------------------
  Ockham-entailment is stronger than super-entailment: whenever $A$
  Ockham-entails $B$, then $A$ super-entails $B$, but not the other way around.

  Suppose $A$ Ockham-entails $B$. Let $t$ be any time in any temporal model at
  which $A$ is true, i.e.: true relative to all histories through $t$. Since $A$
  Ockham-entails $B$, $B$ is true at $t$ relative to all histories through
  $t$. So $A$ super-entails $B$.

  But suppose $A$ super-entails $B$. Let $t$ be any time on any history $h$ in
  any temporal model at which $A$ is true. We can't infer that $B$ is true at
  $t$ on $h$, for $A$ may be false at $t$ relative to other histories $h'$. So
  we can't infer that $A$ Ockham-entails $B$. Indeed, $\tF p$ super-entails
  $\Box \tF p$, but $\tF p$ does not Ockham-entail $\Box \tF p$.

  \subsection*{Exercise 7.18}

  % ------------------------------------------------------------------------
% file `ml-exercise-98-solution-body.tex'
%
%     solution of type `exercise' with id `98'
%
% generated by the `solution' environment of the
%   `xsim' package v0.16a (2020/01/16)
% from source `ml' on 2020/09/26 on line 1472
% ------------------------------------------------------------------------
  $\tN A \to A$ is valid. So $\Box(\tN A \to A)$. By assumption (a),
  we could infer $\tG\Box(\tN A \to A)$. Assumption (b) would tell us
  that $\tG(\Box(\tN A \to A) \to (\tN A \to A))$. By the \textbf{FK}
  principle and \emph{modus ponens}, this implies
  $\tG\Box(\tN A \to A) \to \tG(\tN A \to A)$. So if we make both
  assumptions, we can infer $\tG(\tN A \to A)$, which is equivalent to
  the claim that whatever is now true is always going to be true.

  It is not obvious which of the two assumptions should be rejected.

  \subsection*{Exercise 8.1}

  % ------------------------------------------------------------------------
% file `ml-exercise-99-solution-body.tex'
%
%     solution of type `exercise' with id `99'
%
% generated by the `solution' environment of the
%   `xsim' package v0.16a (2020/01/16)
% from source `ml' on 2020/09/26 on line 281
% ------------------------------------------------------------------------
  For example: $\neg A \strictif A$ or $(B\lor \neg B) \strictif A$.

  \subsection*{Exercise 8.2}

  % ------------------------------------------------------------------------
% file `ml-exercise-100-solution-body.tex'
%
%     solution of type `exercise' with id `100'
%
% generated by the `solution' environment of the
%   `xsim' package v0.16a (2020/01/16)
% from source `ml' on 2020/09/26 on line 309
% ------------------------------------------------------------------------
  $W = \{ w \}$, $R = \emptyset$, $V(p,w)=1$, $V(q,w)=0$.

  \subsection*{Exercise 8.5}

  % ------------------------------------------------------------------------
% file `ml-exercise-103-solution-body.tex'
%
%     solution of type `exercise' with id `103'
%
% generated by the `solution' environment of the
%   `xsim' package v0.16a (2020/01/16)
% from source `ml' on 2020/09/26 on line 510
% ------------------------------------------------------------------------
  All the examples from section 1 also work in the subjunctive mood. For
  \textbf{P4} and \textbf{P5} this yields:
  \begin{itemize}
  \item If our opponents had been cheating, we would never have found out.
    Therefore: If we had found out that our opponents are cheating, then
    they wouldn't have been cheating.
  \item If you had added sugar to your coffee, it would have tasted
    good. Therefore: If you had added sugar and vinegar to your coffee, it would
    have tasted good.
  \end{itemize}

  \subsection*{Exercise 8.6}

  % ------------------------------------------------------------------------
% file `ml-exercise-104-solution-body.tex'
%
%     solution of type `exercise' with id `104'
%
% generated by the `solution' environment of the
%   `xsim' package v0.16a (2020/01/16)
% from source `ml' on 2020/09/26 on line 565
% ------------------------------------------------------------------------
  Suppose $A \to B$ is assertable. Then $A\to B$ is known. So $\Kn (A \to
  B)$. In S4, it follows that $\Kn\Kn(A \to B)$. So the epistemically strict
  conditional $\Kn(A \to B)$ is assertable. Conversely, if $\Kn(A \to B)$ is
  assertable, then it is known; so $\Kn\Kn(A \to B)$. In S4, it follows that
  $\Kn(A \to B)$. So $A \to B$ is assertable.


  \subsection*{Exercise 8.8}

  % ------------------------------------------------------------------------
% file `ml-exercise-106-solution-body.tex'
%
%     solution of type `exercise' with id `106'
%
% generated by the `solution' environment of the
%   `xsim' package v0.16a (2020/01/16)
% from source `ml' on 2020/09/26 on line 811
% ------------------------------------------------------------------------
  Suppose $A \boxright B$ is true at some world $w$ in some model $M$. So $B$ is
  true at all the closest $A$-worlds to $w$. Now either $A$ is true at $w$ or
  $A$ is false at $w$. If $A$ is false at $w$, then $A\to B$ is true at $w$. If
  $A$ is true at $w$, then $w$ is one of the closest $A$-worlds to $w$, by Weak
  Centring; so $B$ is true at $w$; and so $A\to B$ is true at $w$. Either way,
  then, $A\to B$ is true at $w$.


  \subsection*{Exercise 8.9}

  % ------------------------------------------------------------------------
% file `ml-exercise-107-solution-body.tex'
%
%     solution of type `exercise' with id `107'
%
% generated by the `solution' environment of the
%   `xsim' package v0.16a (2020/01/16)
% from source `ml' on 2020/09/26 on line 824
% ------------------------------------------------------------------------
  If $A$ is true at no worlds, then $\emph{Min}^{\prec_w}(\{u: M,u\models A\})$
  is the empty set. So it is vacuously true that $M,v \models B$ for all
  $v \in \emph{Min}^{\prec_w}(\{ u: M,u \models A \})$.

  \subsection*{Exercise 8.13}

  % ------------------------------------------------------------------------
% file `ml-exercise-111-solution-body.tex'
%
%     solution of type `exercise' with id `111'
%
% generated by the `solution' environment of the
%   `xsim' package v0.16a (2020/01/16)
% from source `ml' on 2020/09/26 on line 1300
% ------------------------------------------------------------------------
  `All dogs are barking': $\forall x(Dx \to Bx)$\\
  `Some dogs are barking': $\exists x(Dx \land Bx)$\\
  `Most dogs are barking' cannot be translated in terms of $\Mostly x$. We need
  a binary quantifier: $\Mostly x(Bx / Dx)$

  \subsection*{Exercise 9.1}

  % ------------------------------------------------------------------------
% file `ml-exercise-112-solution-body.tex'
%
%     solution of type `exercise' with id `112'
%
% generated by the `solution' environment of the
%   `xsim' package v0.16a (2020/01/16)
% from source `ml' on 2020/09/26 on line 109
% ------------------------------------------------------------------------
  \begin{enumerate}[(a)]
  \item $Srj \land Szj$; \qquad $r$: Keren, $z$: Keziah, $j$: Jemima, $S$: -- is a sister of --
  \item $\forall x (Mx \to Ox)$; \qquad $M$: -- is a myriapod, $O$: -- is oviparous
  \item $\exists x Dx$; \qquad $D$: -- is at the door
  \item $\neg\forall x (Sx \to Lxl)$; \qquad $l$: logic; $S$: -- is a student, $L$: -- loves --
  \item $\forall x (Sx \land Lxl \to \exists y Lxy)$; \qquad $l$: logic; $S$: -- is a student, $L$: -- loves --
  \end{enumerate}


  \subsection*{Exercise 9.3}

  % ------------------------------------------------------------------------
% file `ml-exercise-114-solution-body.tex'
%
%     solution of type `exercise' with id `114'
%
% generated by the `solution' environment of the
%   `xsim' package v0.16a (2020/01/16)
% from source `ml' on 2020/09/26 on line 314
% ------------------------------------------------------------------------
  For both cases, use $Fx$ as sentence $A$, and $\neg Fx$ as $B$.

  \subsection*{Exercise 9.4}

  % ------------------------------------------------------------------------
% file `ml-exercise-115-solution-body.tex'
%
%     solution of type `exercise' with id `115'
%
% generated by the `solution' environment of the
%   `xsim' package v0.16a (2020/01/16)
% from source `ml' on 2020/09/26 on line 468
% ------------------------------------------------------------------------
  Let $M_1$ be a model with two worlds; each world can see the other but not
  itself; at both worlds, all sentence letters are false. Let $M_2$ be a model
  with a single world that can see itself and at which all sentence letters are
  false. The very same $\L_M$-sentences are true at all worlds in these models
  (as a simple proof by induction shows). But the first model is irreflexive and
  the second isn't.

  \subsection*{Exercise 9.5}

  % ------------------------------------------------------------------------
% file `ml-exercise-116-solution-body.tex'
%
%     solution of type `exercise' with id `116'
%
% generated by the `solution' environment of the
%   `xsim' package v0.16a (2020/01/16)
% from source `ml' on 2020/09/26 on line 481
% ------------------------------------------------------------------------
  For example: $\forall w Rww$ states that $R$ is reflexive; as the answer to
  the previous exercise shows, this can't be expressed in $\L_M$. There are many
  other examples.

  \subsection*{Exercise 9.7}

  % ------------------------------------------------------------------------
% file `ml-exercise-118-solution-body.tex'
%
%     solution of type `exercise' with id `118'
%
% generated by the `solution' environment of the
%   `xsim' package v0.16a (2020/01/16)
% from source `ml' on 2020/09/26 on line 650
% ------------------------------------------------------------------------
  If a sentence is valid in first-order predicate logic, then a fully expanded
  tree for the sentence will close and show that the sentence is valid. But if a
  sentence is not valid, the tree might grow forever. There is no algorithm for
  detecting whether a tree will grow forever. So there are trees that grow
  forever but you can't tell that they will grow forever.

  \subsection*{Exercise 9.8}

  % ------------------------------------------------------------------------
% file `ml-exercise-119-solution-body.tex'
%
%     solution of type `exercise' with id `119'
%
% generated by the `solution' environment of the
%   `xsim' package v0.16a (2020/01/16)
% from source `ml' on 2020/09/26 on line 799
% ------------------------------------------------------------------------

  \begin{enumerate}[(a)]
  \item $\Box Fa$ \\ $a$: John, $F$: -- is hungry.\\[-2mm]
  \item $\Box \forall x(Fx \to Gx)$ \\ $F$: -- is a cyclist, $G$: -- has legs.\\[1mm]
    Also correct (but different in meaning): $\forall x (Fx \to \Box Gx)$.\\
    Close but incorrect: $\forall x \Box(Fx \to Gx)$.\\[-2mm]
  \item $\forall x (Fx \to \Diamond Gx)$ \\ $F$: -- is a day, $G$: -- is our last day.\\[1mm]
    The English sentence could also be understood as
    $\Diamond \forall x (Fx \to Gx)$, but that would be a very strange
    thing to say.\\[-2mm]
  \item $a=b \to \Diamond \forall x(Fx \to Hxb)$ \\ $a$: water, $b$: H$_2$O, $F$: -- is a cucumber, $H$: -- contains --.\\[1mm]
    It's not obvious that `water' and `H$_2$O' function as names. So you might also try to translate them as predicates:\\[1mm]
    $\Box\forall x (Gx \leftrightarrow Jx) \to \Diamond \forall x (Fx \to
    \exists y(Jy \land Hxy))$\\
    $F$: -- is a cucumber, $H$: -- contains --, $G$: -- is water, $J$: -- is H$_2$O.\\[1mm]
    Here $\Box(Gx \leftrightarrow Jx)$ is an attempt to say that water is H$_2$O.
    \\[-2mm]


  \item %If anyone wants to leave early, they should do so quietly.
    $\forall x \Ob(Fx \to Gx)$\\
    $F$: -- wants to leave early, $G$: -- leaves quietly.\\[1mm]
    Even better, if we can use the conditional obligation operator: $\forall x \Ob(Gx / Fx)$\\[1mm]
    %
    These aren't too far off either:
    $\forall x (Fx \to \Ob Gx)$, $\Ob\forall x(Fx \to Gx)$
    \\[-2mm]

  \item % Everyone who bought a ticket is allowed to enter.

    $\forall x (\exists y (Fy \land Hxy) \to \Pe Gx)$\\
    $F$: -- is a ticket, $G$: -- enters, $H$: -- bought --.\\[1mm]
    Perhaps even better: $\forall x \Pe(Gx/ \exists y (Fy \land Hxy))$.\\[1mm]
    You could translate `bought a ticket' as a simple predicate here;
    you could also use a temporal operator to account for the past
    tense of `bought' (but it's confusing to use two different kinds
    of `$\tP$' in one sentence).
    \\[-2mm]

  %   % Delete this one:
  % \item Alice knows that there is someone who knows that she has a lover.
  % \item This sentence can't be translated with our resources. We would
  %   have to change the knowledge operator $\Kn$ into a two-place
  %   operator taking a subject and a proposition as argument. Then we
  %   could say:\\[1mm]
  %   $\Kn(a, \exists x \Kn(x, \exists y Hyx))$\\
  %   $a$: Alice, $Hxy$: $x$ is $y$'s lover.
  %   \\[-2mm]

  \item $\tF \forall x (\tN Fx \to Gx)$\\
    $F$: -- is rich, $G$: -- is poor.\\[1mm]
    The non-trivial reading of the sentence can't be translated without the
    `now' operator $\tN$.

  \end{enumerate}

  \subsection*{Exercise 9.11}

% ------------------------------------------------------------------------
% file `ml-exercise-122-solution-body.tex'
%
%     solution of type `exercise' with id `122'
%
% generated by the `solution' environment of the
%   `xsim' package v0.16a (2020/01/16)
% from source `ml' on 2020/09/26 on line 1004
% ------------------------------------------------------------------------
  \tree{
    \nnode{15}{1.}{$b=c$}{}{}\\
    \nnode{15}{2.}{$A$}{}{}\\                 % e.g. Fc v Hcb, want to get Fb v Hcb
    \nnode{15}{3.}{$b=b$}{}{(Id.)}\\
    \nnode{15}{4.}{$c=b$}{}{\hspace{32mm}(1, 3, LL (first version))}\\
    \nnode{15}{5.}{$A[b//c]$}{}{\hspace{32mm}(4, 2, LL (first version))}
  }

  \subsection*{Exercise 9.12}

% ------------------------------------------------------------------------
% file `ml-exercise-123-solution-body.tex'
%
%     solution of type `exercise' with id `123'
%
% generated by the `solution' environment of the
%   `xsim' package v0.16a (2020/01/16)
% from source `ml' on 2020/09/26 on line 1072
% ------------------------------------------------------------------------
  \begin{enumerate}[(a)]
  \item $\exists x \exists y(Fx \land Fy \land x\not=y \land \forall z(Fz \to (z=x \lor z=y)))$
  \item $\forall x \forall y\forall z\forall v(Fx \land Fy \land Fz \land Fv \to (x=y \lor x=z \lor x=v \lor y=z \lor y=v \lor z=v))$
  \end{enumerate}


  \subsection*{Exercise 9.13}

  % ------------------------------------------------------------------------
% file `ml-exercise-124-solution-body.tex'
%
%     solution of type `exercise' with id `124'
%
% generated by the `solution' environment of the
%   `xsim' package v0.16a (2020/01/16)
% from source `ml' on 2020/09/26 on line 1150
% ------------------------------------------------------------------------
  The \emph{de dicto} reading of (a) can be translated as
  \begin{equation*}
    \Diamond \exists x (Px \land \forall y(Py \to x\!=\!y) \land x\!=\!c),%
  \end{equation*}
  where `$P$' translates `-- is 45th US President' and `$c$' denotes Hilary
  Clinton. The \emph{de re} reading can be translated as
  \begin{equation*}
    \exists x (Px \land \forall y(Py \to x\!=\!y) \Diamond x\!=\!c).
  \end{equation*}

  \subsection*{Exercise 10.1}

  % ------------------------------------------------------------------------
% file `ml-exercise-125-solution-body.tex'
%
%     solution of type `exercise' with id `125'
%
% generated by the `solution' environment of the
%   `xsim' package v0.16a (2020/01/16)
% from source `ml' on 2020/09/26 on line 164
% ------------------------------------------------------------------------
  (b), (c), (d), and (e) are true; (a), (f), (g) are false.

  \subsection*{Exercise 10.3}

  % ------------------------------------------------------------------------
% file `ml-exercise-127-solution-body.tex'
%
%     solution of type `exercise' with id `127'
%
% generated by the `solution' environment of the
%   `xsim' package v0.16a (2020/01/16)
% from source `ml' on 2020/09/26 on line 309
% ------------------------------------------------------------------------
  \begin{enumerate}[(a)]
  \item $W=\{ w \}$, $wRw$, $D = \{ \text{Alice} \}$, $V(F,w) = \{ \text{Alice} \}$, $V(G,w) = \emptyset$
  \item $W=\{ w,v \}$, $wRw$ and $wRv$, $D = \{ \text{Alice}, \text{Bob} \}$, $V(F,w) = \{ \text{Alice} \}$, $V(F,v) = \{ \text{Bob} \}$
  \item $W=\{ w,v \}$, $wRw$ and $wRv$, $D = \{ \text{Alice}, \text{Bob} \}$, $V(F,w) = \{ \text{Alice} \}$, $V(F,v) = \emptyset$
  \item $W=\{ w,v \}$, $wRw$ and $wRv$, $D = \{ \text{Alice}, \text{Bob} \}$, $V(P,w) = \{ \text{Alice} \}$, $V(P,v) = \emptyset$, $V(Q,w) = \emptyset$, $V(Q,v) = \emptyset$
  \end{enumerate}

  \subsection*{Exercise 10.4}

  % ------------------------------------------------------------------------
% file `ml-exercise-128-solution-body.tex'
%
%     solution of type `exercise' with id `128'
%
% generated by the `solution' environment of the
%   `xsim' package v0.16a (2020/01/16)
% from source `ml' on 2020/09/26 on line 531
% ------------------------------------------------------------------------
  $\Box \forall x \exists y(x\!=\!y) \to \forall x \Box\exists y(x\!=\!y)$ is an
  instance of the Converse Barcan Formula. If we read the box as a relevant kind
  of circumstantial necessity, and Loafy could have failed to exist, the
  consequent of this conditional is false. But the antecedent is true.

  \subsection*{Exercise 10.5}

  % ------------------------------------------------------------------------
% file `ml-exercise-129-solution-body.tex'
%
%     solution of type `exercise' with id `129'
%
% generated by the `solution' environment of the
%   `xsim' package v0.16a (2020/01/16)
% from source `ml' on 2020/09/26 on line 557
% ------------------------------------------------------------------------
  (1) is equivalent to the Barcan Formula, (4) to the Converse Barcan
  Formula. (2) is highly implausible. (1) and (4) are often regarded as
  implausible, for the reasons I discuss in the text. (3) is as plausible or
  implausible as the Converse Barcan Formula. (Consider the instance
  $\exists x \Box \neg \exists y(x\!=\!y) \to \Box\exists x \neg \exists
  y(x\!=\!y)$.

  \subsection*{Exercise 10.9}

  % ------------------------------------------------------------------------
% file `ml-exercise-133-solution-body.tex'
%
%     solution of type `exercise' with id `133'
%
% generated by the `solution' environment of the
%   `xsim' package v0.16a (2020/01/16)
% from source `ml' on 2020/09/26 on line 1508
% ------------------------------------------------------------------------
  We would assume that (i) the name $g$ picks out a statue at all accessible
  worlds, (ii) $l$ picks out a lump of clay at all accessible worlds, and (iii)
  at the actual world, $l$ and $g$ pick out the same thing: the statue-shaped
  lump on the shelf.

  \subsection*{Exercise 10.10}

  % ------------------------------------------------------------------------
% file `ml-exercise-134-solution-body.tex'
%
%     solution of type `exercise' with id `134'
%
% generated by the `solution' environment of the
%   `xsim' package v0.16a (2020/01/16)
% from source `ml' on 2020/09/26 on line 1524
% ------------------------------------------------------------------------
  To render $\forall x \forall y (x\!=\!y \to \Box x\!=\!y)$ valid, we can
  stipulate that for any variables $x,y$, and worlds $w,v$ such that $wRv$, and
  any assignment function $g$, if $g(x,w) = g(y,w)$ then $g(x,v)=g(y,v)$. (We
  should make an analogous stipulation for the interpretation function $V$ with
  respect to names.)

  \subsection*{Exercise 10.11}

% ------------------------------------------------------------------------
% file `ml-exercise-135-solution-body.tex'
%
%     solution of type `exercise' with id `135'
%
% generated by the `solution' environment of the
%   `xsim' package v0.16a (2020/01/16)
% from source `ml' on 2020/09/26 on line 1646
% ------------------------------------------------------------------------
  Translation: $\exists x (Tx \land Wx \land \neg K Wx \land \neg K \neg Wx)$,
  where $T$ translates `-- is a ticket' and `-- will win'.

  If variables are directly referential, then this sentence is true in any
  scenario in which I don't know which ticket will win.

  In counterpart semantics, we could assume that the sentence makes \emph{being
    the winning ticket} a salient respect of similarity, so that a ticket $t$ at
  an epistemically accessible world $w$ is a counterpart of the winning ticket
  at the actual world iff $t$ is winning at $w$. This renders the sentence false
  in any situation in which I know that there is a winning ticket (as the
  first part of the sentence conveys).

\end{document}
